\begin{figure}[H]\centering
\begin{tikzpicture}[x=1cm,y=1cm]
\draw[figline,very thick] (0,0)--(10,0);
\draw[figgray] (0,-0.12)--(-0.3,-0.6)--(0.3,-0.6)--cycle;
\draw[figgray] (10,-0.12)--(9.7,-0.6)--(10.3,-0.6)--cycle;
\draw[figaxis,figred] (2,1.5)--(2,0.15);
\draw[figaxis,figred] (6,1.5)--(6,0.15);
\node[above,font=\small] at (2,1.5) {carga 1};
\node[above,font=\small] at (6,1.5) {carga 2};
\node[below,font=\small,align=center] at (5,-0.85) {misma matriz $A$; distintos vectores $\mathbf b$\\una factorización, varias resoluciones};
\end{tikzpicture}
\caption{Esquema del puente usado para motivar la reutilización de LU. Las flechas representan cargas posibles, sin magnitudes ni geometría estructural especificadas por la transcripción. Cambiar las cargas modifica el lado derecho del sistema.}
\end{figure}
